% s-047-p047-m001; original PDF 47, 48
\Aphorism{s-047-p047-m001}{38}
Staaten sind doch nur „so zu sagen“ Organismen, und ein nicht ganz mythologischer Mensch wird die oft durchgeführte Parallele zwischen Thier- und Gesellschaftskörper kaum für mehr halten als für eine geistreiche Allegorie. Aber schon, dass sich überhaupt eine solche Parallele ziehen lässt, stimmt zu wohlwollendem Mitleid gegen den Staat — vorausgesetzt nämlich, dass der Anblick des Menschen zu wohlwollendem Mitleid stimmt. Hegel begriff es anders: ihm war der Staat die vollkommenste aller Abstractionen, ein Absolutum ohne Anfang, Wachsthum und Ende, er spricht mit frommen Schaudern von dieser höchsten Incarnation der Idee und meint, es sei „herb“ den Staat zu fassen. Diese Nüance der Ehrfurcht, einer politomanen Zeit und einem politomanen Gehirne entsprungen, dürften wir uns heute gründlich abgewöhnt haben; wir denken über den Staat utilitarisch, wie es der Streit der Interessen und Interessengruppen uns gelehrt hat. Wohlverstanden, ich zweifle nicht, dass es heute noch begeisterte Monarchisten oder Republicaner giebt, aber diese Art Enthusiasmus gilt dem ästhetischen Reiz geschichtlich überlieferter Gebilde; eine Idolatrie vor der Staatsform als solcher, eine Anbetung der reinen Gesetzlichkeit als der höchsten Verwirklichung sittlicher Ideale ist nirgends mehr vorhanden, mag sie noch so oft geheuchelt werden. Uns ist der Staat, mögen wir uns conservativ oder subversiv gebärden, ein Nebeneinander widerstreitender Egoismen, ein Etagenbau von Functionen und Organen, und insofern wirklich ein ungefähres Gleichniss des menschlichen Leibes. Haben wir also Geduld mit dem Staate, wie wir Geduld mit unserem Leibe haben! Wundern wir uns nicht, wenn im Gebaren dieses Riesenkörpers Sinn und Vernunft manchmal zu fehlen scheinen; auch in unserem Leibesleben ist ja das Gehirn nicht immer oberste Instanz. Handeln, denken, fühlen wir nicht lange Lebensstrecken hin unter dem Druck von Ganglien, Magen, Eingeweiden und Geschlechtstheilen, haben wir nicht unsere Zeiten der Unfruchtbarkeit, ja des Rückschrittes, verschlafen und verschlemmen wir nicht zwei Drittel unseres Schaffenstages, — wie? und der Staat sollte geistiger, cerebraler, philosophischer leben als wir, die Staatsangehörigen? dem Riesenkörper sollten die Widerlichkeiten, die schlimmen Zufälle und Irrationalitäten erspart sein, denen alles Körperliche ausgesetzt ist?
\par

